\chapter{Introduction}
\label{ch:intro}

\section{What this workspace is}

\texttt{baxter\_ros2\_jazzy} runs a Rethink Robotics Baxter Research Robot from
ROS~2 Jazzy. The default path is simulation: Gazebo Harmonic, \texttt{ros2\_control}
arm controllers, and MoveIt~2 planning against them. On top of that sits a
hardware path that talks to the robot's own ROS~1 Noetic master through a bridge,
and that path has moved a real arm --- under supervision, in dated sessions.

Baxter shipped in 2012 with a ROS~1 Groovy/Indigo SDK. Rethink Robotics closed in
2018 and there is no vendor ROS~2 support, so anything running Baxter on Jazzy is
reconstructing the command path from the robot's wire protocol outward. That is
what this workspace does, and it is why the hardware chapters read as carefully as
they do: nothing here is a supported vendor path.

\section{How to read the support claims}

Every capability statement in this manual falls into one of three categories, and
the manual is written so you can always tell which:

\begin{itemize}
  \item \textbf{Passed} --- a gate ran and its evidence is recorded. Simulation
    gates run in CI on every push.
  \item \textbf{Passed, supervised} --- it worked on one physical robot, on a named
    date, with a person holding the e-stop. This is not general support.
  \item \textbf{Deferred} or \textbf{not measured} --- said plainly rather than
    left to inference.
\end{itemize}

\begin{safetybox}[The evidence base, stated once]
Every hardware claim in this manual comes from \emph{one} Baxter, driven under
supervision, on four dated occasions: 2026-07-22 (bridge and interlocks),
2026-07-24 (first supervised motion, first MoveIt execution), and 2026-07-25
(speed characterisation and a MoveIt re-run). Sustained duty, gripper commands,
and any second robot are unmeasured. Do not read this manual as evidence that
Baxter is generally supported on ROS~2.
\end{safetybox}

The canonical machine-readable version of this is
\texttt{docs/support\_matrix.md} in the repository; if that file and this manual
ever disagree, the file wins, because CI checks it and nothing checks prose.

\section{What works}

\begin{table}[h]
\centering
\small
\begin{tabularx}{\linewidth}{@{}lllX@{}}
\toprule
\textbf{Profile} & \textbf{Status} & \textbf{Chapter} & \textbf{Notes} \\
\midrule
Gazebo sim & passed & \ref{ch:sim} & Both arms, 17 joint states, reversible motion. \\
Gazebo + RViz & passed & \ref{ch:sim} & RobotModel/TF profile. \\
MoveIt~2 sim & passed & \ref{ch:moveit} & OMPL, single- and dual-arm groups. \\
MoveIt RViz & passed & \ref{ch:moveit} & MotionPlanning panel, IK markers. \\
Hardware bridge & supervised & \ref{ch:bridge} & 2026-07-22. \\
Supervised motion & supervised & \ref{ch:hwmotion} & 2026-07-24; speed limits 2026-07-25. \\
MoveIt on hardware & supervised, \emph{limited} & \ref{ch:hwmotion} & Left arm aborts most goals at default speed. \\
Grippers, cameras, Zenoh & deferred & --- & No gate requires them yet. \\
\bottomrule
\end{tabularx}
\end{table}

\section{What this manual will not tell you}

It does not contain raw safety-topic publishing or hardware enable commands.
Those exist in the repository's operational runbook, behind procedures that
assume a supervised session and an e-stop within reach. A manual that teaches
someone to energise a 2-metre dual-arm robot from a code snippet is a manual that
has misjudged its audience.

It also does not document the Rethink Intera manufacturing workflow --- task
training, vacuum grippers, Modbus, the factory pick-and-place tooling. That is a
different product built on the same hardware. If you have a manufacturing Baxter
and want to train pick-and-place tasks by demonstration, this workspace is not
what you want.

\section{Layout}

Chapter~\ref{ch:platform} covers the robot itself: joints, limits, and the
concepts the rest of the manual assumes. Chapters~\ref{ch:install}
to~\ref{ch:moveit} are the simulation path, which needs no robot and is where
anyone should start. Chapters~\ref{ch:bridge} and~\ref{ch:hwmotion} are the
hardware path. Chapter~\ref{ch:verify} is how to check that a change did not break
anything, without a robot. The appendices are the condensed command sheet and the
numeric tables.
